\documentclass[10pt,aspectratio=169]{beamer}

% Paquetes
\usepackage[utf8]{inputenc}
\usepackage[T1]{fontenc}
\usepackage[spanish]{babel}
\usepackage{csquotes}
\usepackage{graphicx}
\usepackage{xcolor}
\usepackage{booktabs}
\usepackage{array}
\usepackage{ragged2e}
\usepackage{tikz}
\usetikzlibrary{positioning, arrows.meta, shapes, fit, backgrounds}

% Tema y colores
\usetheme{Madrid}
\usecolortheme{whale}

% Colores personalizados (paleta azul similar al informe)
\definecolor{azulPrincipal}{RGB}{31,73,125}
\definecolor{azulClaro}{RGB}{68,114,196}
\definecolor{azulSuave}{RGB}{220,230,241}
\definecolor{grisOscuro}{RGB}{64,64,64}
\definecolor{acento}{RGB}{237,125,49}
\definecolor{verdeExito}{RGB}{84,130,53}

\setbeamercolor{frametitle}{fg=white, bg=azulPrincipal}
\setbeamercolor{title}{fg=white, bg=azulPrincipal}
\setbeamercolor{structure}{fg=azulPrincipal}
\setbeamercolor{block title}{fg=white, bg=azulClaro}
\setbeamercolor{block body}{fg=grisOscuro, bg=azulSuave!40}
\setbeamercolor{alertblock title}{fg=white, bg=acento}
\setbeamercolor{alertblock body}{fg=grisOscuro, bg=acento!15}
\setbeamercolor{normal text}{fg=grisOscuro}
\setbeamercolor{itemize item}{fg=azulPrincipal}
\setbeamercolor{itemize subitem}{fg=azulClaro}
\setbeamercolor{enumerate item}{fg=azulPrincipal}

% Quitar iconos de navegación
\setbeamertemplate{navigation symbols}{}
\setbeamertemplate{footline}[frame number]
\setbeamertemplate{headline}{}
\setbeamertemplate{itemize item}[circle]
\setbeamertemplate{itemize subitem}[square]
\setbeamertemplate{blocks}[rounded][shadow=false]

% Información del documento
\title{\textbf{Large Language Models en\\Sistemas de Información Empresariales}}
\subtitle{Integración de LLM en ERP, CRM, SCM, BI, gestión documental y atención al cliente}
\author{Grupo 2}
\institute{\small Integrantes: Martínez Fuentes, K. \textperiodcentered\ Vergara, J. \textperiodcentered\ Motta, J. \textperiodcentered\ Triveño, F. \textperiodcentered\ Guerrero, C. \textperiodcentered\ Sernaqué, J.}
\date{Curso: Tendencias en Sistemas de Información\\
Docente: Angulo Calderón, César Augusto\\
Lima, Perú -- 2026}

\begin{document}

% =============== DIAPOSITIVA PORTADA ===============
\begin{frame}[plain]
    \begin{tikzpicture}[remember picture, overlay]
        \fill[azulPrincipal] (current page.north west) rectangle ([yshift=-5cm]current page.north east);
    \end{tikzpicture}
    \vspace{1cm}
    \begin{center}
        {\color{azulPrincipal}\rule{\textwidth}{2pt}}\\[0.3cm]
        {\Large\textbf{Grupo 2}}\\[1cm]
        {\Huge\textcolor{azulPrincipal}{\textbf{Large Language Models}}}\\[0.3cm]
        {\Huge\textcolor{azulPrincipal}{\textbf{en Sistemas de Información}}}\\
        {\Huge\textcolor{azulPrincipal}{\textbf{Empresariales}}}\\[1cm]
        {\large\itshape Integración de LLM en ERP, CRM, SCM, BI,\\gestión documental y atención al cliente}\\[1.5cm]
        \begin{flushleft}
            \textbf{\textcolor{azulPrincipal}{Integrantes:}}\\[0.2cm]
            Martínez Fuentes, Kenny Uldarico Danny \\
            Vergara Pachas, José Luis \\
            Motta Chumbe, JeanPierre D'Angelo \\
            Triveño Daza, Felipe Santiago \\
            Guerrero Falla, Christian Angelo \\
            Sernaqué Gutierrez, José Roberto
        \end{flushleft}
    \end{center}
    \vfill
    \begin{flushright}
        \small
        \textcolor{azulPrincipal}{\textbf{Curso:}} Tendencias en Sistemas de Información\\
        \textcolor{azulPrincipal}{\textbf{Docente:}} Angulo Calderón, César Augusto\\[0.3cm]
        \textbf{Lima, Perú -- 2026}
    \end{flushright}
\end{frame}

% =============== AGENDA ===============
\begin{frame}{Agenda}
    \begin{columns}[T]
        \begin{column}{0.5\textwidth}
            \textbf{\textcolor{azulPrincipal}{1. Introducción y problema}}\\
            \textit{(3 min)}\\[0.4cm]
            \textbf{\textcolor{azulPrincipal}{2. Marco conceptual}}\\
            \textit{(2 min)}\\[0.4cm]
            \textbf{\textcolor{azulPrincipal}{3. Evolución de los SI}}\\
            \textit{(2 min)}\\[0.4cm]
            \textbf{\textcolor{azulPrincipal}{4. Integración con ERP/CRM/SCM/BI}}\\
            \textit{(3 min)}
        \end{column}
        \begin{column}{0.5\textwidth}
            \textbf{\textcolor{azulPrincipal}{5. Caso de uso: Vodafone Ireland}}\\
            \textit{(4 min)}\\[0.4cm]
            \textbf{\textcolor{azulPrincipal}{6. Arquitectura tecnológica}}\\
            \textit{(3 min)}\\[0.4cm]
            \textbf{\textcolor{azulPrincipal}{7. Riesgos y gobernanza}}\\
            \textit{(2 min)}\\[0.4cm]
            \textbf{\textcolor{azulPrincipal}{8. Conclusiones}}\\
            \textit{(1 min)}
        \end{column}
    \end{columns}
    \vspace{0.5cm}
    \begin{block}{\textbf{Duración total: 20 minutos}}
        \centering Trabajo grupal -- Grupo 2\\
        \textit{Tendencias en Sistemas de Información}
    \end{block}
\end{frame}

% =============== INTRODUCCIÓN Y PROBLEMA ===============
\section{Introducción y problema}
\begin{frame}{¿Qué son los LLM?}
    \begin{columns}[T]
        \begin{column}{0.55\textwidth}
            \begin{block}{Definición}
                Modelos de IA basados en la arquitectura de \textbf{transformadores}, entrenados sobre corpus de texto masivo (cientos de miles de millones de tokens) para \textbf{comprender, generar y manipular} lenguaje natural.
            \end{block}
            \vspace{0.3cm}
            \textbf{Modelos representativos:}
            \begin{itemize}
                \item \textbf{GPT-4 / GPT-4o} -- OpenAI
                \item \textbf{Claude 3.5 / 4} -- Anthropic
                \item \textbf{Gemini 1.5 / 2} -- Google
                \item \textbf{LLaMA 3} -- Meta
                \item \textbf{Mistral} -- open weights
            \end{itemize}
        \end{column}
        \begin{column}{0.45\textwidth}
            \begin{alertblock}{Capacidades clave}
                \begin{itemize}
                    \item Generación de texto
                    \item Resumen y traducción
                    \item Razonamiento
                    \item Respuesta a preguntas
                    \item Generación de código
                    \item Análisis de sentimiento
                \end{itemize}
            \end{alertblock}
            \vspace{0.2cm}
            \begin{center}
                \textcolor{azulPrincipal}{\textit{``La nueva generación de aplicaciones empresariales''}}
            \end{center}
        \end{column}
    \end{columns}
\end{frame}

\begin{frame}{Problema central}
    \begin{center}
        \fbox{\parbox{0.85\textwidth}{\centering\vspace{0.3cm}
        \textit{¿Cómo pueden los Large Language Models (LLM) contribuir a mejorar la capacidad analítica, la automatización de flujos de trabajo y la toma de decisiones en las organizaciones modernas mediante su integración en los Sistemas de Información Empresariales?}\vspace{0.3cm}}}
    \end{center}
    \vspace{0.3cm}
    \textbf{\textcolor{azulPrincipal}{Situación actual:}}
    \begin{itemize}
        \item Exceso de información \textbf{no estructurada} (>80\% de los datos empresariales)
        \item Sistemas tradicionales (ERP, CRM, BI) \textbf{no comprenden} lenguaje natural
        \item Procesos manuales lentos, baja automatización cognitiva
        \item Personalización limitada y tiempos de respuesta elevados
    \end{itemize}
\end{frame}

\begin{frame}{Justificación: ¿por qué importa?}
    \begin{columns}[T]
        \begin{column}{0.33\textwidth}
            \begin{block}{Tecnológica}
                \begin{itemize}
                    \item Transformers + self-attention
                    \item Fine-tuning por dominio
                    \item APIs y orquestación
                    \item RAG sobre datos propios
                \end{itemize}
            \end{block}
        \end{column}
        \begin{column}{0.33\textwidth}
            \begin{block}{Organizacional}
                \begin{itemize}
                    \item Optimiza procesos operativos
                    \item Reduce análisis manual
                    \item Mejora toma de decisiones
                    \item Aumenta productividad
                \end{itemize}
            \end{block}
        \end{column}
        \begin{column}{0.33\textwidth}
            \begin{block}{Estratégica}
                \begin{itemize}
                    \item Ventaja competitiva
                    \item Nuevos servicios digitales
                    \item SI proactivos y cognitivos
                    \item Transformación digital
                \end{itemize}
            \end{block}
        \end{column}
    \end{columns}
\end{frame}

% =============== MARCO CONCEPTUAL ===============
\section{Marco conceptual}
\begin{frame}{Conceptos relacionados}
    \begin{columns}[T]
        \begin{column}{0.5\textwidth}
            \textbf{\textcolor{azulPrincipal}{IA Generativa (GenAI)}}\\
            Modelos que crean contenido nuevo (texto, imágenes, audio, código).\\[0.3cm]
            \textbf{\textcolor{azulPrincipal}{RAG (Retrieval-Augmented Generation)}}\\
            Combina un LLM con un sistema de recuperación vectorial para fundamentar respuestas en conocimiento propio de la empresa.\\[0.3cm]
        \end{column}
        \begin{column}{0.5\textwidth}
            \textbf{\textcolor{azulPrincipal}{Embeddings}}\\
            Representaciones vectoriales densas; insumo fundamental de búsqueda semántica.\\[0.3cm]
            \textbf{\textcolor{azulPrincipal}{Agentes Inteligentes}}\\
            LLM como ``cerebro'' que planifica, invoca herramientas y ejecuta acciones.\\[0.3cm]
            \textbf{\textcolor{azulPrincipal}{APIs e Integración}}\\
            LangChain, LlamaIndex, Azure OpenAI, AWS Bedrock, Vertex AI.
        \end{column}
    \end{columns}
\end{frame}

\begin{frame}{LLM vs. sistemas tradicionales}
    \begin{center}
    \small
    \begin{tabular}{p{3.3cm} p{4.8cm} p{4.8cm}}
        \toprule
        \textbf{Dimensión} & \textbf{ERP / CRM / BI} & \textbf{LLM} \\
        \midrule
        Tipo de datos & Estructurados (tablas, registros) & No estructurados (texto, voz, imágenes) \\
        Interacción & Formularios, SQL, dashboards & Lenguaje natural, conversación, voz \\
        Respuesta & Determinística, basada en reglas & Probabilística, generativa \\
        Personalización & Segmentación rígida & Dinámica por usuario y contexto \\
        Adaptabilidad & Configuración / desarrollo & Re-ajuste vía prompting y RAG \\
        Conocimiento & Esquemas y código & Corpus masivo + RAG \\
        Interfaz & GUI rígida, capacitación previa & Conversacional, baja curva \\
        Riesgos & Errores de proceso & Alucinaciones, sesgos, fugas \\
        \bottomrule
    \end{tabular}
    \end{center}
    \vspace{0.3cm}
    \begin{alertblock}{Nivel de madurez}
        \centering Inversión privada en GenAI en 2023: USD \textbf{25.2 mil millones} (8x vs. 2022) -- Stanford HAI (2024). Faltan \textbf{evaluaciones estandarizadas} de responsabilidad.
    \end{alertblock}
\end{frame}

% =============== EVOLUCIÓN ===============
\section{Evolución de los SI}
\begin{frame}{De los sistemas transaccionales a los sistemas autónomos}
    \begin{center}
    \begin{tikzpicture}[node distance=0.4cm, every node/.style={font=\small}]
        \node[draw, rounded corners, fill=azulSuave!60, text width=2.2cm, align=center, minimum height=1.1cm] (t1) {\textbf{1960--90}\\Transaccional\\(ERP, CRM, SCM)};
        \node[draw, -Latex, thick, right=of t1, yshift=-0.3cm] {};
        \node[draw, rounded corners, fill=azulSuave!80, text width=2.2cm, align=center, minimum height=1.1cm, right=1.3cm of t1] (t2) {\textbf{1990--2010}\\Analítico\\(BI, OLAP, DW)};
        \node[draw, rounded corners, fill=azulSuave!90, text width=2.2cm, align=center, minimum height=1.1cm, right=1.3cm of t2] (t3) {\textbf{2010--20}\\Inteligente\\(ML, Data Science)};
        \node[draw, rounded corners, fill=azulClaro, text=white, text width=2.2cm, align=center, minimum height=1.1cm, right=1.3cm of t3] (t4) {\textbf{2020--}\\Generativo\\(ChatGPT, RAG)};
        \node[draw, rounded corners, fill=acento, text=white, text width=2.2cm, align=center, minimum height=1.1cm, right=1.3cm of t4] (t5) {\textbf{Futuro}\\Autónomo\\(AI Agents)};
        \draw[-Latex, thick, azulPrincipal] (t1) -- (t2);
        \draw[-Latex, thick, azulPrincipal] (t2) -- (t3);
        \draw[-Latex, thick, azulPrincipal] (t3) -- (t4);
        \draw[-Latex, thick, acento] (t4) -- (t5);
    \end{tikzpicture}
    \end{center}
    \vspace{0.4cm}
    \begin{block}{Salto cualitativo}
        De predecir / clasificar (ML clásico) a \textbf{crear contenido y asistir cognitivamente} (LLM). Próxima frontera: agentes que \textbf{ejecutan tareas} sin intervención humana.
    \end{block}
\end{frame}

% =============== INTEGRACIÓN ===============
\section{Integración con SI}
\begin{frame}{Integración con ERP y CRM}
    \begin{columns}[T]
        \begin{column}{0.5\textwidth}
            \begin{block}{ERP (SAP, Oracle, Dynamics, Odoo)}
                \begin{itemize}
                    \item Asistentes en lenguaje natural
                    \item Reportes financieros automáticos
                    \item Clasificación de facturas/pedidos
                    \item Soporte a usuarios
                \end{itemize}
                \textit{``¿Cuál fue el margen bruto del Q3 en la región sur?''}
            \end{block}
        \end{column}
        \begin{column}{0.5\textwidth}
            \begin{block}{CRM (Salesforce, HubSpot, Zoho)}
                \begin{itemize}
                    \item Atención 24/7 con chatbots
                    \item Análisis de sentimiento
                    \item Borradores personalizados
                    \item Resumen automático de reuniones
                \end{itemize}
                \vspace{0.2cm}
                \textbf{Vodafone Ireland:} 3x más rápido en crear journeys, 80\% menos utterances malinterpretadas.
            \end{block}
        \end{column}
    \end{columns}
\end{frame}

\begin{frame}{Integración con SCM, BI, gestión documental y mesa de ayuda}
    \begin{columns}[T]
        \begin{column}{0.5\textwidth}
            \textbf{\textcolor{azulPrincipal}{SCM}} -- Cadena de suministro
            \begin{itemize}
                \item Análisis de noticias y disrupciones
                \item Planes de demanda explicables
                \item Asistentes para proveedores (RFP)
            \end{itemize}
            \textbf{\textcolor{azulPrincipal}{BI}} -- Analítica aumentada
            \begin{itemize}
                \item Consultas en lenguaje natural
                \item Generación automática de insights
                \item Storytelling de datos
            \end{itemize}
        \end{column}
        \begin{column}{0.5\textwidth}
            \textbf{\textcolor{azulPrincipal}{Gestión documental}}
            \begin{itemize}
                \item Búsqueda semántica
                \item Resumen y extracción de entidades
                \item RAG con citas verificables
            \end{itemize}
            \textbf{\textcolor{azulPrincipal}{Mesa de ayuda}}
            \begin{itemize}
                \item Clasificación automática de tickets
                \item Resolución autónoma de nivel 1
                \item Generación de artículos de conocimiento
            \end{itemize}
        \end{column}
    \end{columns}
\end{frame}

% =============== CASO DE USO ===============
\section{Caso de uso}
\begin{frame}{Caso: Vodafone Ireland + IBM watsonx Assistant}
    \begin{columns}[T]
        \begin{column}{0.55\textwidth}
            \textbf{Contexto:}\\
            TOBi, el asistente virtual de Vodafone Ireland, atendía a millones de clientes con un modelo anterior. La tasa de resolución estaba \textbf{estancada en 35\%}.\\[0.3cm]
            \textbf{Solución:}\\
            Migración al nuevo IBM watsonx Assistant con \textbf{capacidades generativas} (RAG, NLU basado en foundation model, aprendizaje continuo).\\[0.3cm]
            \textbf{Reconstrucción completa} de los journeys conversacionales y ajustes en la aplicación interna.
        \end{column}
        \begin{column}{0.45\textwidth}
            \begin{alertblock}{Resultados medibles}
                \begin{center}
                {\Huge\color{acento}\textbf{3x}}\\[0.2cm]
                \small más rápido creando nuevos journeys\\[0.4cm]
                {\Huge\color{acento}\textbf{+20\%}}\\[0.2cm]
                \small mejora en tasas de contención\\[0.4cm]
                {\Huge\color{verdeExito}\textbf{--80\%}}\\[0.2cm]
                \small utterances malinterpretadas
                \end{center}
            \end{alertblock}
        \end{column}
    \end{columns}
\end{frame}

% =============== ARQUITECTURA ===============
\section{Arquitectura tecnológica}
\begin{frame}{Arquitectura de 5 capas para integrar LLM}
    \begin{center}
    \begin{tikzpicture}[
        node distance=0.25cm,
        capa/.style={draw, rounded corners=4pt, text width=10cm, align=center, minimum height=1cm, font=\small\bfseries},
        capa1/.style={capa, fill=azulSuave!60},
        capa2/.style={capa, fill=azulSuave!75},
        capa3/.style={capa, fill=azulSuave!85, text=white},
        capa4/.style={capa, fill=azulClaro, text=white},
        capa5/.style={capa, fill=acento, text=white},
        every node/.style={font=\small}
    ]
        \node[capa1] (c5) {5. \textbf{Gobierno}: SSO, OAuth, RBAC/ABAC, auditoría, evaluación, cumplimiento};
        \node[capa2, above=of c5] (c4) {4. \textbf{Aplicación}: copilotos, agentes, chatbots, APIs para sistemas consumidores};
        \node[capa3, above=of c4, fill=azulPrincipal] (c3) {3. \textbf{Modelo}: LLM comercial o self-hosted + RAG + fine-tuning + prompt engineering};
        \node[capa2, above=of c3] (c2) {2. \textbf{Orquestación e indexación}: LangChain, LlamaIndex, bases vectoriales, ETL};
        \node[capa1, above=of c2] (c1) {1. \textbf{Datos}: ERP, CRM, SCM, BI, gestión documental, base de conocimiento};
        \draw[-Latex, thick, azulPrincipal] (c1) -- (c2);
        \draw[-Latex, thick, azulPrincipal] (c2) -- (c3);
        \draw[-Latex, thick, azulPrincipal] (c3) -- (c4);
        \draw[-Latex, thick, acento] (c4) -- (c5);
    \end{tikzpicture}
    \end{center}
    \vspace{0.3cm}
    \begin{block}{Patrón RAG como eje}
        \centering El usuario consulta $\rightarrow$ retrieval semántico $\rightarrow$ contexto al LLM $\rightarrow$ respuesta fundamentada con citas $\rightarrow$ ejecución sobre SI vía APIs
    \end{block}
\end{frame}

% =============== APLICACIONES Y BENEFICIOS ===============
\section{Aplicaciones y beneficios}
\begin{frame}{Aplicaciones y beneficios esperados}
    \begin{columns}[T]
        \begin{column}{0.5\textwidth}
            \textbf{\textcolor{azulPrincipal}{Aplicaciones en SI}}
            \begin{itemize}
                \item \textbf{ERP}: consultas NL, reportes automáticos
                \item \textbf{CRM}: atención 24/7, personalización
                \item \textbf{BI}: analítica aumentada
                \item \textbf{Gestión documental}: búsqueda semántica, RAG
                \item \textbf{Mesa de ayuda}: clasificación, soporte autónomo
            \end{itemize}
        \end{column}
        \begin{column}{0.5\textwidth}
            \textbf{\textcolor{azulPrincipal}{Beneficios}}
            \begin{itemize}
                \item \textbf{Estratégicos}: innovación, ventaja competitiva
                \item \textbf{Operativos}: menos tiempos, más eficiencia
                \item \textbf{Tecnológicos}: APIs flexibles, escalabilidad
                \item \textbf{Analíticos}: decisiones basadas en datos no estructurados
                \item \textbf{Económicos}: reducción de costes, ROI
                \item \textbf{Humanos}: mejor UX, mayor satisfacción
            \end{itemize}
        \end{column}
    \end{columns}
\end{frame}

% =============== RIESGOS ===============
\section{Riesgos y gobernanza}
\begin{frame}{Riesgos clave y estrategias de mitigación}
    \small
    \begin{center}
    \begin{tabular}{p{3.5cm} p{5cm} p{5cm}}
        \toprule
        \textbf{Riesgo} & \textbf{Descripción} & \textbf{Mitigación} \\
        \midrule
        Alucinaciones & Respuestas inventadas & RAG + validación humana \\
        Fuga de información & Datos sensibles en prompts & RBAC, cifrado, modelos privados \\
        Sesgos & Respuestas discriminatorias & Auditoría, datasets balanceados \\
        Prompt injection & Manipulación maliciosa & Filtros, separación instrucciones/datos \\
        Dependencia & Vendor lock-in & Arquitectura multi-modelo \\
        Costos & Tokens, cómputo, talento & Optimización, SLM, caching \\
        Cumplimiento & RGPD, AI Act & Gobierno de IA, DPIA \\
        \bottomrule
    \end{tabular}
    \end{center}
    \vspace{0.3cm}
    \begin{alertblock}{Marcos de referencia}
        \centering \textbf{NIST AI RMF} (gestión de riesgos) \textperiodcentered\ \textbf{NIST AI 600-1 Generative AI Profile} \textperiodcentered\ \textbf{OWASP Top 10 for LLM Applications 2025}
    \end{alertblock}
\end{frame}

% =============== TENDENCIAS ===============
\section{Tendencias}
\begin{frame}{Tendencias futuras (3--5 años)}
    \begin{columns}[T]
        \begin{column}{0.5\textwidth}
            \begin{block}{Modelo y arquitectura}
                \begin{itemize}
                    \item \textbf{LLM multimodales} (texto, imagen, audio)
                    \item \textbf{Small Language Models} (SLM) en el edge
                    \item \textbf{IA privada} on-premise por industria
                \end{itemize}
            \end{block}
            \begin{block}{Aplicación}
                \begin{itemize}
                    \item \textbf{Agentes autónomos} multi-equipo
                    \item Automatización cognitiva de procesos (RPA + LLM)
                    \item Copilotos nativos en SAP, Salesforce, ServiceNow
                \end{itemize}
            \end{block}
        \end{column}
        \begin{column}{0.5\textwidth}
            \begin{block}{Gobernanza}
                \begin{itemize}
                    \item \textbf{AI Act} europeo + ISO/IEC 42001
                    \item Trazabilidad, explicabilidad obligatoria
                    \item Evaluaciones de impacto
                \end{itemize}
            \end{block}
            \begin{alertblock}{Visión prospectiva}
                \centering SI \textbf{multimodales, autónomos, auditables y energéticamente eficientes}.\\[0.2cm]
                \small Los LLM se integrarán de forma nativa en los ERPs, CRMs, BI y mesas de ayuda, dando lugar a la próxima generación de Sistemas de Información.
            \end{alertblock}
        \end{column}
    \end{columns}
\end{frame}

% =============== CONCLUSIONES ===============
\section{Conclusiones}
\begin{frame}{Conclusiones}
    \begin{columns}[T]
        \begin{column}{0.5\textwidth}
            \begin{block}{Tecnológica}
                Los LLM mejoran la interacción persona--sistema mediante interfaces conversacionales, búsqueda semántica y automatización cognitiva, sustentadas en transformers, RAG y agentes.
            \end{block}
            \begin{block}{Organizacional}
                Su adopción transforma finanzas, ventas, soporte, logística, RR.\,HH. y atención al cliente. El caso Vodafone Ireland confirma mejoras medibles (3x, +20\%, --80\%).
            \end{block}
        \end{column}
        \begin{column}{0.5\textwidth}
            \begin{alertblock}{Crítica}
                Se requiere gestionar alucinaciones, fugas, prompt injection, sesgos, dependencia, costos y cumplimiento (NIST AI RMF, NIST GenAI Profile, OWASP).
            \end{alertblock}
            \begin{block}{Prospectiva}
                SI multimodales, agentes autónomos, modelos privados, IA regulada y energéticamente eficiente. El LLM será una \textbf{capa nativa} de los SI empresariales.
            \end{block}
        \end{column}
    \end{columns}
\end{frame}

% =============== CIERRE ===============
\begin{frame}[plain]
    \begin{tikzpicture}[remember picture, overlay]
        \fill[azulPrincipal] (current page.south west) rectangle ([yshift=4cm]current page.south east);
    \end{tikzpicture}
    \vfill
    \begin{center}
        {\Huge\textcolor{azulPrincipal}{\textbf{¿Preguntas?}}}\\[1.5cm]
        {\Large\textcolor{azulPrincipal}{\textbf{Gracias por su atención}}}\\[1cm]
        \begin{flushright}
            \small
            \textcolor{azulPrincipal}{\textbf{Grupo 2}}\\
            Martínez Fuentes, Vergara, Motta,\\
            Triveño, Guerrero, Sernaqué\\[0.3cm]
            \textit{Curso: Tendencias en Sistemas de Información}
        \end{flushright}
    \end{center}
    \vfill
    \begin{tikzpicture}[remember picture, overlay]
        \node[anchor=south west, xshift=1cm, yshift=0.5cm] at (current page.south west) {\color{white}\small\textbf{Lima, Perú -- 2026}};
    \end{tikzpicture}
\end{frame}

% =============== APÉNDICE: REFERENCIAS ===============
\appendix
\begin{frame}{Referencias bibliográficas}
    \begin{enumerate}\setlength{\itemsep}{0pt}
        \item \footnotesize \textbf{NIST} (2023). \textit{Artificial Intelligence Risk Management Framework (AI RMF 1.0)}. DOI: 10.6028/NIST.AI.100-1.
        \item \textbf{IBM} (2024). \textit{Vodafone Ireland -- Redesigning conversational assistant to improve performance}. IBM Case Studies.
        \item \textbf{NIST} (2024). \textit{AI RMF Generative Artificial Intelligence Profile}. DOI: 10.6028/NIST.AI.600-1.
        \item \textbf{Salesforce} (2024). \textit{State of the AI Connected Customer}.
        \item \textbf{Stanford HAI} (2024). \textit{The 2024 AI Index Report: Measuring trends in AI}.
        \item \textbf{OWASP} (2024). \textit{Top 10 for LLM Applications 2025}. OWASP Foundation.
        \item \footnotesize Cao, S. et al. (2026). \textit{Reframing Conversational Design in HRI: Deliberate Design with AI Scaffolds}. arXiv:2601.12084.
        \item \footnotesize He, J. \& Liu, J. (2026). \textit{Seeing to Think? How Source Transparency Design Shapes Interactive Information Seeking and Evaluation in Conversational AI}. arXiv:2601.14611.
        \item \footnotesize Vijayakumar, H. (2026). \textit{UX in the Age of AI: Rethinking Evaluation Metrics Through a Statistical Lens}. arXiv:2605.05600.
    \end{enumerate}
\end{frame}

\end{document}
